% GENERATED FILE. Edit canonical lessons, then run npm run generate:books.
% canonical-source-hash: fnv1a64:baf0797bf00d0524
% canonical-lessons: ZH-C197-wenlu, ZH-C197-zhaodao, ZH-W197-zhao-seek, ZH-C197-milu, ZH-C197-yizhi, ZH-C197-luxian

\chapter{Finding the Way}
\label{ch:zh-finding-the-way}
\hlchaptermodality{197}

\begin{chapteropening}
\textbf{By the end of this chapter:} \emph{I can ask the way, say I have found a place or got lost, and go straight on.}

Complete the last lesson of chapter 197: I can ask the way, say I have found a place or got lost, and go straight on.
\end{chapteropening}

\section[\texorpdfstring{\zh{问路} (wènlù)}{问路 (wènlù)}]{\zh{问路} (wènlù) — to ask the way}

\label{lesson:ZH-C197-wenlu}



\begin{quote}
Before the new one: say the Mandarin for to take off, then the Mandarin for to reach a destination.
\end{quote}

\subsection*{You'll want to know: \zh{问路}}
\textbf{\zh{问路}} — \emph{wènlù} — \textquotedblleft{}to ask the way\textquotedblright{}.

\begin{quote}
\textbf{\zh{他在问路}} \emph{tā zài wènlù} — He's asking the way.
\end{quote}

\begin{cousinweb}
\textbf{\zh{问路}} (\emph{wènlù}) is built from \zh{问} 'to ask' and \zh{路} 'road', both inherited from Middle Chinese, \emph{mjunH} and \emph{luH} (Guangyun, in Baxter's transcription).
\end{cousinweb}

\begin{grammarlens}[title={What you've built}]
One more word for asking where things are.
\end{grammarlens}

\subsection*{Your turn}
\begin{itemize}
\raggedright
  \item \emph{Say it:} \emph{wènlù}
  \item \emph{Say it:} \emph{wènlù}, once more
  \item \emph{Say it:} \emph{tā zài wènlù} — He's asking the way.
  \item \emph{Recall:} say the Mandarin for to take off, then the Mandarin for to reach a destination, then say \emph{wènlù} again
\end{itemize}

\begin{tcolorbox}[breakable,colback=teal!4,colframe=teal!35!black,title={Before you move on}]
What does \zh{问路} mean? (\textquotedblleft{}To ask the way\textquotedblright{}.) Say it once more.
\end{tcolorbox}
\section[\texorpdfstring{\zh{找到} (zhǎodào)}{找到 (zhǎodào)}]{\zh{找到} (zhǎodào) — to find}

\label{lesson:ZH-C197-zhaodao}



\begin{quote}
Before the new one: say the Mandarin for to reach a destination, then the Mandarin for to ask the way.
\end{quote}

\subsection*{You'll want to know: \zh{找到}}
\textbf{\zh{找到}} — \emph{zhǎodào} — \textquotedblleft{}to find\textquotedblright{}.

\begin{quote}
\textbf{\zh{我找到书店了}} \emph{wǒ zhǎodào shūdiàn le} — I've found the bookshop.
\end{quote}

The new character, in pieces: \zh{找} — seven strokes: on the left the hand piece of \zh{打}; on the right a bar, a long slanting hook, a stroke falling to the left across it, and a dot. The next lesson puts it in your hand.

\begin{cousinweb}
\textbf{\zh{找到}} (\emph{zhǎodào}) is built from \zh{找} 'to look for' and \zh{到} 'to arrive': looking that gets there. \zh{到} is inherited from Middle Chinese \emph{tawH} (Guangyun, in Baxter's transcription); no dictionary we use records an older form of \zh{找}, so its history is not settled.
\end{cousinweb}

\begin{grammarlens}[title={What you've built}]
One more word for asking where things are.
\end{grammarlens}

\subsection*{Your turn}
\begin{itemize}
\raggedright
  \item \emph{Say it:} \emph{zhǎodào}
  \item \emph{Say it:} \emph{zhǎodào}, once more
  \item \emph{Say it:} \emph{wǒ zhǎodào shūdiàn le} — I've found the bookshop.
  \item \emph{Recall:} say the Mandarin for to reach a destination, then the Mandarin for to ask the way, then say \emph{zhǎodào} again
\end{itemize}

\begin{tcolorbox}[breakable,colback=teal!4,colframe=teal!35!black,title={Before you move on}]
What does \zh{找到} mean? (\textquotedblleft{}To find\textquotedblright{}.) Say it once more.
\end{tcolorbox}
\section[\texorpdfstring{\zh{找} (zhǎo)}{找 (zhǎo)}]{\zh{找} — the hand piece, then a hook and a fall}

\label{lesson:ZH-W197-zhao-seek}



\begin{quote}
The word first, then the character inside it.

\begin{itemize}
\raggedright
  \item \emph{Recall:} say \emph{zhǎodào} — \textbf{R1}, one lesson back, before it decays
\end{itemize}

\begin{quote}
\zh{打} \zh{钱}
\end{quote}

\emph{Read it:} each one aloud, with nothing said first — \emph{dǎ}, \emph{qián}
\end{quote}

\subsection*{Script — the hand piece, then a hook and a fall}
\hlblockfigure{\detokenize{figures/ZH-W197-zhao-seek-filmstrip.pdf}}{How \zh{找} is written, stroke by stroke}

\begin{quote}
\zh{找}
\end{quote}

\textbf{Two pieces, left before right.}

\textbf{The left piece, the hand piece, three strokes,} the same three as in \zh{打}: a short bar, a vertical that hooks to the left at its foot, then a stroke rising to the right across it.

\textbf{The right piece, four strokes.} A bar. A long slanting hook from the top down to the lower right, flicking up at its foot. A stroke falling to the left across it. Then a dot at the top right.

3 + 4 = seven. The right piece is the right piece of \zh{钱} with one bar fewer.

\subsection*{Your turn}
\begin{itemize}
\raggedright
  \item \emph{Write it:} \zh{找}, all seven, with the model in front of you
  \item \emph{Cover:} write it again from memory, and count the pieces, not the strokes
\end{itemize}

\begin{tcolorbox}[breakable,colback=teal!4,colframe=teal!35!black,title={Before you move on}]
How many strokes in \zh{找}? (\textbf{Seven}.) Where has your hand made the left piece before? (\textbf{In \zh{打}}.) What comes last? (\textbf{The dot at the top right}.)

Source: \href{https://raw.githubusercontent.com/chanind/hanzi-writer-data/68d10a4b21150cae5e1ebbd223eed289cf32d90c/data/\%E6\%89\%BE.json}{Hanzi Writer Data \zh{找}, pinned stroke paths and medians}.
\end{tcolorbox}
\section[mílù]{mílù — to lose your way; to get lost}

\label{lesson:ZH-C197-milu}



\begin{quote}
Before the new one: say the Mandarin for to ask the way, then the Mandarin for to find.
\end{quote}

\subsection*{You'll want to know: mílù}
\emph{mílù} — \textquotedblleft{}to lose your way; to get lost\textquotedblright{}.

\begin{quote}
\emph{wǒ mílù le} — I'm lost.
\end{quote}

Its characters come later; for now it is a word for the ear.

\begin{cousinweb}
\emph{mílù} is built from \emph{mí} 'confused' and \zh{路} 'road', both inherited from Middle Chinese, \emph{mej} and \emph{luH} (Guangyun, in Baxter's transcription), which glosses the first 'to be bewildered'.
\end{cousinweb}

\begin{grammarlens}[title={What you've built}]
One more word for asking where things are.
\end{grammarlens}

\subsection*{Your turn}
\begin{itemize}
\raggedright
  \item \emph{Say it:} \emph{mílù}
  \item \emph{Say it:} \emph{mílù}, once more
  \item \emph{Say it:} \emph{wǒ mílù le} — I'm lost.
  \item \emph{Recall:} say the Mandarin for to ask the way, then the Mandarin for to find, then say \emph{mílù} again
\end{itemize}

\begin{tcolorbox}[breakable,colback=teal!4,colframe=teal!35!black,title={Before you move on}]
What does mílù mean? (\textquotedblleft{}To lose your way; to get lost\textquotedblright{}.) Say it once more.
\end{tcolorbox}
\section[yìzhí]{yìzhí — straight on; all the way}

\label{lesson:ZH-C197-yizhi}



\begin{quote}
Before the new one: say the Mandarin for to find, then the Mandarin for to lose your way.
\end{quote}

\subsection*{You'll want to know: yìzhí}
\emph{yìzhí} — \textquotedblleft{}straight on; all the way\textquotedblright{}.

\begin{quote}
\emph{cóng zhèr yìzhí zǒu} — Go straight on from here.
\end{quote}

Its characters come later; for now it is a word for the ear.

\zh{一} changes its tone here: yì, before the rising tone of zhí.

\begin{cousinweb}
\emph{yìzhí} is built from \zh{一} 'one' and \emph{zhí} 'straight', both inherited from Middle Chinese, \emph{'jit} and \emph{drik} (Guangyun, in Baxter's transcription), which glosses the second 'upright'.
\end{cousinweb}

\begin{grammarlens}[title={What you've built}]
One more word for asking where things are.
\end{grammarlens}

\subsection*{Your turn}
\begin{itemize}
\raggedright
  \item \emph{Say it:} \emph{yìzhí}
  \item \emph{Say it:} \emph{yìzhí}, once more
  \item \emph{Say it:} \emph{cóng zhèr yìzhí zǒu} — Go straight on from here.
  \item \emph{Recall:} say the Mandarin for to find, then the Mandarin for to lose your way, then say \emph{yìzhí} again
\end{itemize}

\begin{tcolorbox}[breakable,colback=teal!4,colframe=teal!35!black,title={Before you move on}]
What does yìzhí mean? (\textquotedblleft{}Straight on; all the way\textquotedblright{}.) Say it once more.
\end{tcolorbox}
\section[lùxiàn]{lùxiàn — route}

\label{lesson:ZH-C197-luxian}



\begin{quote}
Before the new one: say the Mandarin for to lose your way, then the Mandarin for straight on.
\end{quote}

\subsection*{You'll want to know: lùxiàn}
\emph{lùxiàn} — \textquotedblleft{}route\textquotedblright{}.

\begin{quote}
\emph{qù jīchǎng de lùxiàn hěn jiǎndān} — The route to the airport is simple.
\end{quote}

Its characters come later; for now it is a word for the ear.

\begin{cousinweb}
\emph{lùxiàn} is built from \zh{路} 'road' and \emph{xiàn} 'a line', both inherited from Middle Chinese, \emph{luH} and \emph{sjenH} (Guangyun, in Baxter's transcription), which glosses the second 'thread'.
\end{cousinweb}

\begin{grammarlens}[title={What you've built}]
One more word for asking where things are.
\end{grammarlens}

\subsection*{Your turn}
\begin{itemize}
\raggedright
  \item \emph{Say it:} \emph{lùxiàn}
  \item \emph{Say it:} \emph{lùxiàn}, once more
  \item \emph{Say it:} \emph{qù jīchǎng de lùxiàn hěn jiǎndān} — The route to the airport is simple.
  \item \emph{Recall:} say the Mandarin for to lose your way, then the Mandarin for straight on, then say \emph{lùxiàn} again
\end{itemize}

\begin{tcolorbox}[breakable,colback=teal!4,colframe=teal!35!black,title={Before you move on}]
What does lùxiàn mean? (\textquotedblleft{}Route\textquotedblright{}.) Say it once more.
\end{tcolorbox}
